\chapter[Introdução]{Introdução}
\label{cap_1}
% ---
\section{Caracterização do Problema}
% ---

A segmentação semântica constitui uma tarefa fundamental em Visão Computacional, com aplicações em áreas como análise de imagens médicas, veículos autônomos, monitoramento por vídeo e realidade aumentada. Embora os avanços recentes impulsionados pelo Aprendizado Profundo (\ac{DL}) tenham ampliado significativamente o desempenho dos modelos, persiste um obstáculo crítico: a necessidade de grandes volumes de imagens anotadas manualmente. Esse processo é custoso, demorado e sujeito a inconsistências, pois requer que especialistas rotulem pixel a pixel cada instância da base de dados.

Outro desafio recorrente diz respeito ao desequilíbrio entre classes, comum em bases reais, o que compromete a capacidade de generalização dos modelos. Esse problema torna-se ainda mais evidente em cenários de alta resolução (como imagens 4K contendo múltiplos objetos), em que tanto o esforço humano quanto o custo computacional se tornam significativamente mais elevados.

Diante desse contexto, o problema central desta pesquisa pode ser formulado da seguinte forma: como propor e validar estratégias que integrem Aprendizado Ativo, Aprendizado por Reforço, superpixels e Aprendizado Contínuo, de modo a reduzir o esforço de anotação em imagens de alta resolução, assegurando qualidade e eficiência computacional no processo de segmentação semântica?

% ---
\section{Motivações e Desafios}
% ---

A principal motivação deste trabalho é reduzir o custo e o tempo necessários para a anotação de imagens, mantendo ou até mesmo aprimorando a qualidade da segmentação semântica gerada pelos modelos. A literatura apresenta diversas estratégias que podem contribuir para esse objetivo, tais como:

\begin{itemize}

\item Aprendizado Ativo (\ac{AL}): busca diminuir a quantidade de amostras que precisam ser anotadas, priorizando aquelas mais informativas para o aprendizado. Embora promissora, a maioria dos métodos ainda depende de heurísticas manuais e opera em ciclos offline, sem interação em tempo real com especialistas humanos \cite{kim2023adaptive,konyushkova2017learning}. %【kim2023adaptive; konyushkova2017learning】.

\item Integração de \ac{AL} com \ac{RL}: abordagens recentes exploram políticas de seleção de dados guiadas por recompensas de desempenho, aumentando a eficiência da escolha de amostras \cite{fang2017active}. %fang2017active.
Entretanto, esses modelos ainda carecem de mecanismos de interação contínua com especialistas.

\item Segmentação por superpixels: possibilita que a anotação seja realizada em unidades visuais semanticamente coerentes, o que reduz ambiguidades de fronteira e esforço de anotação. Segundo Vargas et al. (2014), essa representação permite descritores de imagem em nível mais alto quando comparada à abordagem pixel a pixel.

\item Esquecimento catastrófico em Aprendizado Contínuo: modelos incrementais frequentemente perdem conhecimento adquirido ao serem expostos a novos dados. Estratégias como a Destilação de Conhecimento surgem como alternativas para mitigar esse problema e garantir estabilidade no refinamento do modelo.

\end{itemize}

Diante desse cenário, esta pesquisa propõe avançar o estado da arte por meio de um framework integrado, que combine \ac{AL}, \ac{RL}, superpixels e técnicas de Aprendizado Contínuo, promovendo um processo de anotação mais eficiente, interativo e escalável.


%Com o advento do aprendizado de máquina profundo (Deep Leaning), muitas tarefas de visão computacional, incluindo a segmentação semântica, tem evoluído nos últimos anos \cite{kim2023adaptive}. Apesar dos avanços, umas das principais dificuldades ainda permanece: a necessidade de grandes volumes de dados anotados, que demanda um esforço considerável e tempo elevado. Além disso, é comum que os dados apresentem forte desbalanceamento entre as classes, o que afeta negativamente o aprendizado \cite{casanova2020reinforced}.

%Para mitigar esse problema, estratégias de Active Learning (\ac{AL}) têm sido amplamente estudadas. O objetivo central da \ac{AL} é reduzir o número de amostras que precisam ser rotuladas, identificando aquelas que são mais informativas para o aprendizado. A maioria dos métodos atualmente existentes ainda são baseados em heurísticas manuais \cite{konyushkova2017learning}

%Abordagens recentes propõem a integração de \ac{AL} com Aprendizado por reforço (Reinforcement Learning (\ac{RL})), permitindo que o sistema aprenda políticas de obtenção de dados através do sistema de recompensas de desempenho \cite{fang2017active}. Entretanto, modelos assim operam em ciclos de batchs ou offline, não possuindo uma interação em tempo real com os especialista humanos.

%Além disso, modelos baseados em superpixel têm demonstrado serem promissores ao permitir que as unidades anotáveis se baseiem em unidades visuais com maior nível de semântica, o que reduz a ambiguidade de fronteira e o esforço de anotação. A segmentação de imagem em superpixels tem como objetivo agrupar pixels em regiões homogêneas que capturam a redundância da imagem. Segundo \cite{vargas2014superpixel}, "a representação por superpixels permite descritores de imagem de nível mais alto em comparação com a representação por pixel".


\section{Objetivos}

\subsection{Objetivo Geral}

% Este trabalho tem como objetivo principal propor, desenvolver e validar estratégias eficazes para reduzir o custo e o tempo necessário para anotação de imagens e realizar uma interação em tempo real com especialistas para otimizar essas anotações.

Este trabalho tem como objetivo principal propor, desenvolver e validar um framework interativo para segmentação semântica de imagens, capaz de reduzir o custo e o tempo de anotação manual, ao mesmo tempo em que promove interação em tempo real entre especialistas humanos e o modelo de aprendizado profundo. A proposta busca integrar técnicas de \ac{AL}, \ac{CL}, métodos baseados em superpixels aplicados tanto ao aumento de dados quanto à análise estrutural das regiões de interesse para auxiliar o \ac{AL} e abordagens fundamentadas em \ac{RL} destinadas à melhoria progressiva do modelo, de forma a otimizar o processo de anotação, melhorar a qualidade das segmentações e mitigar problemas como o esquecimento catastrófico.

\subsection{Objetivos Específicos}

%Desenvolver objetivo de data augmentation utilizando superpixel-mix com active learning

\begin{itemize}

%\item Levantamento, análise e seleção de bases de imagens com multiclasses e alta qualidade disponíveis.

\item Realizar levantamento e análise crítica de bases de imagens multiclasse de alta qualidade, selecionando aquelas mais adequadas para validação da proposta.

%\item Levantamento, análise e seleção de técnicas de superpixel para redução de tempo computacional

\item Estudar, selecionar e adaptar técnicas de segmentação por superpixels, visando reduzir o custo computacional e o esforço de anotação, mantendo a precisão das fronteiras dos objetos.

%\item Levantamento e proposta de técnicas de \ac{AL} para identificar amostras mais adequadas, dada as incertezas do modelo.

\item Investigar e propor estratégias de \ac{AL} para identificar amostras mais informativas com base na incerteza do modelo, otimizando o processo de refinamento interativo.

%\item Levantamento, análise e seleção de técnicas de \ac{RL} para reduzir o número de interações da \ac{AL} com o especialista, considerando a interação com as amostras que possuam maior grau de incerteza.

\item Explorar e integrar abordagens de \ac{RL} como apoio ao \ac{AL}, a fim de reduzir o número de interações necessárias com o especialista, priorizando instâncias de maior relevância.

%\item Proposta de inclusão de técnicas de critérios no processo de seleção de dados para retornar ao especialista quanto a incerteza em multiclasses.

\item Propor critérios adicionais para seleção de dados em cenários multiclasse, fornecendo feedback ao especialista sobre regiões de maior incerteza e auxiliando na priorização das correções.

%\item Avaliação e comparação extensiva das técnicas propostas utilizando métricas Dice para comparativo com outras técncias já existentes.

\item Implementar e avaliar a integração das técnicas propostas em um framework interativo, realizando comparações extensivas com métodos existentes, utilizando métricas como Dice e \ac{IoU}.

\item Validar a eficácia do framework em cenários práticos de segmentação semântica, verificando sua capacidade de reduzir tempo e custo de anotação, mitigar o esquecimento catastrófico e transferir conhecimento entre instâncias da mesma classe.

\end{itemize}

\section{Contribuições Esperadas}

Espera-se que o presente trabalho traga contribuições relevantes tanto no campo científico quanto no prático. Do ponto de vista metodológico, a principal contribuição será a proposição e validação de um framework interativo para segmentação semântica de imagens em alta resolução, capaz de reduzir o custo e o tempo de anotação por especialistas, ao mesmo tempo em que minimiza a necessidade de interação constante durante o processo de treinamento.

A integração de estratégias de \ac{AL} com \ac{RL}, aliada ao uso de técnicas de segmentação por superpixels, pode representar um avanço em relação às abordagens tradicionais, oferecendo maior eficiência na seleção de amostras e maior precisão na segmentação de imagens 4K contendo múltiplos objetos. Espera-se também obter contribuições a partir de novas estratégias de seleção de amostras mais informativas para as técnicas de aprendizado ativo.

Além disso, espera-se que a utilização de técnicas de \ac{CL} e Destilação de Conhecimento permita mitigar o problema do esquecimento catastrófico, garantindo a retenção do conhecimento previamente adquirido pelo modelo ao longo do refinamento interativo.

No âmbito prático, espera-se contribuir para a obtenção e disponibilização de conjuntos de dados anotados de alta qualidade, gerados a partir das pesquisas e experimentos desenvolvidos neste trabalho, de forma a apoiar futuras investigações e aplicações em visão computacional. Dessa forma, a pesquisa deverá não apenas avançar o estado da arte em métodos de anotação assistida, mas também fornecer ferramentas concretas que possam ser incorporadas em cenários reais de análise de imagens.

%A pesquisa tem como princiapal objetivo a otimização das anotações por especialista, considerando tempo de anotação por imagme de alta qualidade e treinamento da rede neural. Tem-se como foco central, melhorar o tempo e custo computacional de imagens de alta qualidades que possuam diversos objetos para segmentação semântica, sem que haja uma interação constante com o especialista.

%Para alcançar esse objetivo, espera-se obter contribuições com a proposta de novas estratégias de \ac{AL} alidas a \ac{RL} e utilização de superpixel que sejam mais adequadas para imagens 4k.

%Também espera-se contribuir para a obtenção e disponbilização de conjuntos de dados anotados, obtidos através das pesquisas e da aplicação da estratégia nas imagens.


\section{Questões e Hipóteses da Pesquisa}

A presente pesquisa busca responder a um conjunto de questões relacionadas à otimização do processo de anotação de imagens em alta resolução para segmentação semântica, considerando eficiência, qualidade e viabilidade prática. Essas questões estão diretamente alinhadas aos objetivos gerais e específicos do trabalho e são acompanhadas de hipóteses que orientam as etapas experimentais.

\begin{itemize}
\item Q1. Como reduzir o custo e o tempo de anotação de imagens em alta resolução (4K) para segmentação semântica, mantendo a qualidade das anotações realizadas por especialistas?

\begin{itemize}
\item H1. O framework interativo proposto reduzirá significativamente o tempo e o custo de anotação, preservando a qualidade das segmentações obtidas.
\end{itemize}

\item Q2. De que forma estratégias de \ac{AL}, aliadas
ao \ac{RL}, podem ser utilizadas para selecionar amostras mais informativas e reduzir a necessidade de interação constante com o especialista?

\begin{itemize}
\item H2. A integração de \ac{AL} e \ac{RL} permitirá a seleção mais eficiente de amostras, diminuindo o número de interações necessárias com o especialista.
\end{itemize}

\item Q3. Qual é o impacto da utilização de técnicas baseadas em superpixels na eficiência da anotação e no desempenho dos modelos de segmentação em imagens complexas com múltiplos objetos?

\begin{itemize}
\item H3. A segmentação por superpixels reduzirá a complexidade computacional do processo, mantendo a precisão das fronteiras dos objetos em imagens multiclasse.
\end{itemize}

\item Q4. Como incorporar técnicas de \ac{CL} e Destilação de Conhecimento para mitigar o problema do esquecimento catastrófico durante o refinamento interativo do modelo?

\begin{itemize}
\item H4. A combinação de \ac{CL} e Destilação de Conhecimento mitigará o esquecimento catastrófico, permitindo a retenção do conhecimento previamente adquirido pelo modelo.
\end{itemize}

\item Q5. De que maneira o framework proposto pode contribuir para a geração de conjuntos de dados anotados de alta qualidade, aplicáveis em diferentes contextos de visão computacional?

\begin{itemize}
\item H5. O framework possibilitará a criação de conjuntos de dados anotados de alta qualidade, que poderão ser disponibilizados para a comunidade científica e utilizados em diferentes aplicações de visão computacional.
\end{itemize}

\end{itemize}

%\begin{itemize}

%\item Como reduzir o custo e o tempo de anotação de imagens em alta resolução (4K) para segmentação semântica, mantendo a qualidade das anotações realizadas por especialistas?

%\item De que forma estratégias de Aprendizado Ativo (AL), aliadas à Aprendizado por Reforço (RL), podem ser utilizadas para selecionar amostras mais informativas e reduzir a necessidade de interação constante com o especialista?

%\item Qual é o impacto da utilização de técnicas baseadas em superpixels na eficiência da anotação e no desempenho dos modelos de segmentação em imagens complexas com múltiplos objetos?

%\item Como incorporar técnicas de Aprendizado Contínuo (Continual Learning) e Destilação de Conhecimento para mitigar o problema do esquecimento catastrófico durante o refinamento interativo do modelo?

%\item De que maneira o framework proposto pode contribuir para a geração de conjuntos de dados anotados de alta qualidade, aplicáveis em diferentes contextos de visão computacional?

%\end{itemize}


\section{Publicações}

Encontra-se em elaboração um artigo científico para divulgação dos resultados preliminares obtidos a partir de experimentos com estratégias de aumento de dados. A submissão deste artigo está prevista para uma conferência da área no primeiro semestre de 2026.
Os resultados subsequentes, relacionados à aplicação e avaliação das estratégias de \ac{AL}, \ac{RL} e segmentação por superpixels, serão organizados e submetidos para publicação em periódicos de relevância na área, de modo a garantir ampla disseminação científica e impacto acadêmico. Além disso, pretende-se publicar em periódico especializado os conjuntos de dados gerados e utilizados nesta pesquisa, assegurando sua reprodutibilidade e promovendo a ampla divulgação dos resultados junto à comunidade.

%Encontra-se em elaboração um artigo científico, para publicação dos resultados preliminares obtidos, a partir de experimentos relacionados à aplicação e avaliação de algumas estratégias de aumento de dados. A submissão deste artigo para publicação em conferência da área está prevista para o segundo semestre de 2026. Os demais resultados a serem obtidos, envolvendo as estrategias de aprendizado ativo, aprendizado por reforço e superpixels, devem também ser publicados em periódicos relevantes na área.